\section{Revue de littérature}
\label{sec:art}


% La prédiction de l'issue de matchs de football est un problème traité par un grand nombre de personnes, 
% tant par des passionnés de foot, que des parieurs aguerris mais aussi 
% par des statisticiens. 
L'un des premiers résultats majeurs dans ce domaine a été établi par 
M. J. Maher \cite{maher}, lorsqu'il expose le fait que le nombre de buts marqués lors de rencontres 
de football peut être simulé par une loi Poisson prenant comme paramètres la force 
d'attaque et de défense des équipes.

\subsection{Approche statistique classique}

Les travaux de Maher \cite{maher} ont posé les fondements de la modélisation 
du football en supposant que les buts marqués suivent des lois de Poisson 
indépendantes. Mark Dixon et Stuart Coles \cite{dixon} ont prolongé cette approche 
en introduisant deux améliorations majeures, la correction de la dépendance via l'ajout 
d'un facteur multiplicatif $\tau$ pour corriger la sous-représentation des faibles 
scores ($0-0, 1-0, 0-1, 1-1$) dans les modèles de Poisson standards. De plus, ils ajoutent une 
pondération temporelle par l'intégration d'une fonction de pondération décroissante 
$\phi(t-t_k)$ permettant de donner plus d'importance aux résultats récents et 
de suivre la dynamique de performance des équipes.

\noindent Bien que très interprétable grâce aux paramètres d'attaque ($\alpha$) et de 
défense ($\beta$), cette approche reste limitée par l'absence de variables 
contextuelles (blessures, cartons) et une sensibilité au choix du paramètre de pondération.

%%%%%%%%%%%%%%%%%%%%%%%%%%%%%%%%%%%%%%%%
%%%%%%%%%%%%%%%%%%%%%%%%%%%%%%%%%%%%%%%%
\subsubsection{Modèle de Poisson -- Maher (1982)}

L'utilisation de la loi de Poisson repose sur l'hypothèse que la possession du ballon se transforme 
en but avec une probabilité $p$ faible, mais que le nombre d'occasions est élevé. 
On notera que la loi binomiale négative constitue une alternative plus flexible, 
capable de modéliser la surdispersion éventuelle des scores, mais nous retenons 
ici la loi de Poisson conformément à Maher (1982)\cite{maher}.\\

\noindent On note $i$ l'équipe à domicile et $j$ l'équipe extérieure. Le couple de 
scores $(x_{ij}, y_{ij})$ est modélisé comme suit : $X_{ij}$ suit une loi de Poisson 
de moyenne $\alpha_i \beta_j$, $Y_{ij}$ suit une loi de Poisson de moyenne 
$\gamma_i \delta_j$, et $X_{ij}$ et $Y_{ij}$ sont supposés indépendants. 
Les paramètres ont l'interprétation suivante : $\alpha_i$ représente la force 
d'attaque de l'équipe $i$ à domicile, $\beta_j$ la faiblesse défensive de l'équipe 
$j$ à l'extérieur, $\gamma_i$ la faiblesse défensive de $i$ à domicile, et $\delta_j$ 
la force d'attaque de $j$ à l'extérieur.\\

\noindent Pour assurer l'identifiabilité du modèle, on impose les contraintes :
\[
\sum_i \alpha_i = \sum_i \beta_i \qquad \text{et} \qquad \sum_i \gamma_i = \sum_i \delta_i
\]
car les produits $\alpha_i \beta_j$ et $\gamma_i \delta_j$ sont invariants par 
rééchelonnage simultané des deux facteurs.\\

\noindent La log-vraisemblance pour les buts à domicile s'écrit :
\[
\log L(\underline{\alpha}, \underline{\beta}) = \sum_{i} \sum_{j \neq i} 
\left( -\alpha_i \beta_j + x_{ij} \log(\alpha_i \beta_j) - \log(x_{ij}!) \right)
\]

Les conditions du premier ordre donnent :
\[
\frac{\partial \log L}{\partial \alpha_i} = \sum_{j \neq i} 
\left( -\beta_j + \frac{x_{ij}}{\alpha_i} \right) = 0
\]

\noindent ainsi, les estimateurs du maximum de vraisemblance $\hat{\alpha}_i$ et $\hat{\beta}_j$ 
vérifient le système implicite :
\[
\hat{\alpha}_i = \frac{\sum_{j \neq i} x_{ij}}{\sum_{j \neq i} \hat{\beta}_j}
\qquad \text{et} \qquad
\hat{\beta}_j = \frac{\sum_{i \neq j} x_{ij}}{\sum_{i \neq j} \hat{\alpha}_i}
\]

Ce système peut être résolu par Newton-Raphson, ce qui conduit à l'approximation :
\[
\hat{\alpha}_i = \frac{\sum_{j \neq i} x_{ij}}{\sqrt{S_x}}, \qquad
\hat{\beta}_j = \frac{\sum_{i \neq j} x_{ij}}{\sqrt{S_x}}, \qquad
S_x = \sum_i \sum_{j \neq i} x_{ij}
\]
où $S_x$ désigne le nombre total de buts marqués à domicile sur l'ensemble des matchs.
Les paramètres $\hat{\gamma}_i$ et $\hat{\delta}_j$ sont estimés de façon analogue 
à partir des $y_{ij}$.



%%%%%%%%%%%%%%%%%%%%%%%%%%%%%
%%%%%%%%%%%%%%%%%%%%%%%%%%%%%
\subsubsection{Approche statistique de Dixon et Coles}
Dans la continuité du modèle de Maher, Dixon et Coles ont ajouté un coefficient $\gamma$ 
représentant l'avantage de jouer à domicile, ainsi qu'une fonction $\tau$ pour gérer 
les petits scores $x$-$y$ avec $x,y \in \left\{ 0,1 \right\} $. En effet, l'hypothèse 
d'indépendance des buts initiée par Maher est relaxée par Dixon et Coles pour les 
faibles scores mais conservée pour les scores plus importants. Ils ajoutent également 
une fonction $\phi$ pondérant les matchs récents.




%%%%%%%%%%%%%%%%%%%%%%%%%%%%%%%%%%%%%%%%
%%%%%%%%%%%%%%%%%%%%%%%%%%%%%%%%%%%%%%%%

\subsection{Différentes modélisations}

Egidi et Torelli \cite{egidi} enrichissent l'analyse en comparant, via une 
approche bayésienne, deux types de modèles. Dans un premier temps "Result-based" par 
une modélisation multinomiale de l'issue (1, X, 2) utilisant des modèles probit
\footnote{Le modèle probit utilise la fonction de répartition de la loi normale 
centrée réduite pour modéliser la probabilité d'un événement binaire ou multinomial.} 
ou logit\footnote{Le modèle logit transforme les probabilités (comprises entre 0 et 1) 
en valeurs réelles via la fonction log-odds $ln(p/(1-p)$), facilitant ainsi les calculs 
en inférence bayésienne.}. 
Cette approche est particulièrement pertinente lors des phases finales de tournois où 
les écarts de niveau sont faibles, bien qu'elle ignore l'ampleur des scores 
(ex: une victoire 5-0 est traitée comme un 1-0). Puis, "Goal-based", ici la modélisation 
est basée sur une loi de Poisson bivariée permettant d'inclure un paramètre de 
dépendance ($\lambda_3$) entre les buts. Elle s'avère plus efficace lorsque 
les écarts de niveau sont élevés.


\noindent L'utilisation de l'inférence par chaînes de Markov de Monte Carlo (MCMC) permet 
une quantification précise de l'incertitude, bien qu'elle nécessite des ressources 
de calcul importantes par rapport au maximum de vraisemblance.

\subsection{Apprentissage automatique}

Carloni \textit{et al.} \cite{carloni} ainsi que Fátima Rodrigues et Ângelo Pinto \cite{rodrigues} explorent le potentiel des algorithmes 
d'apprentissage supervisé (KNN, SVM, Random Forest, Naive Bayes) et des réseaux 
de neurones artificiels (ANN) à quatre couches. Contrairement aux approches 
précédentes, ces modèles ne nécessitent aucune hypothèse a priori sur la 
distribution des buts. 
Leur méthodologie repose sur un nettoyage rigoureux des données et l'utilisation de la \textit{prediction accuracy} et du 
retour sur investissement (ROI) comme métriques de performance. Bien que l'ANN 
soit identifié comme le modèle le plus performant (avec un ROI de 26\% sur un 
échantillon test), ces modèles souffrent d'un manque d'explicabilité et d'une 
sensibilité au découpage des données (\textit{train/test split}), qui peut ne 
pas respecter la chronologie des matchs.
